% !TEX program = xelatex
% !TEX root = ../../TFM.tex
\documentclass[../../TFM.tex]{subfiles}

\begin{document}

% =====================================================
%  DATOS Y CONSTRUCCIÓN DE VARIABLES
% =====================================================

	\chapter{Datos y construcción de variables}
	\label{chap:datos-variables}
	\thesisstartchap{datos-variables}

	\section{Fuentes de datos y muestra}
	\label{sec:fuentes-datos-muestra}

	Este capítulo no replica el detalle metodológico sobre criterios de selección muestral ni sobre la justificación de cada indicador. Su propósito es presentar una caracterización descriptiva del panel consolidado antes de la estimación econométrica. En consecuencia, aquí se identifican de manera general las fuentes de información empleadas, su cobertura y la lógica de integración de la base, mientras que la discusión completa sobre selección de países, período de análisis y especificación del modelo se desarrolla en el capítulo metodológico.

	De manera sintética, la base combina información de comercio internacional, estructura productiva, desempeño macroeconómico y calidad institucional. El \textit{Economic Complexity Index} ($ECI$), $DRES$, $HHI$, $DIVX$, $PEXP$ y $FEXP$ se obtienen o construyen a partir del \textit{Atlas of Economic Complexity}; $RENTS$, $OILPC$, $GASPC$, $COALPC$, $INNOV$, $NET$, $LOG\_GDPPC$, $GOVCONS$ y $FIN$ utilizan insumos de los \textit{World Development Indicators}; $INST$ emplea tres estimadores de los \textit{Worldwide Governance Indicators}; $VOL$ utiliza el índice CTOT del Fondo Monetario Internacional con ponderaciones fijas; y $HUMCAP$ y $RER$ proceden exclusivamente de Penn World Table 11.0.

	La selección mediante $DRES\geq20\%$ identifica 55 países, pero esta cifra
	no equivale necesariamente al número final de países de cada estimación.
	La captura oficial de $ECI$ mediante la API del Atlas cubre los 55 países y
	los 1.430 país--años de 1996--2021. Los cuatro componentes conjuntos de
	$RENTS$ tienen información para 53 países; $HUMCAP$ para 47; $RER$ para 52;
	$VOL$ para 54; y
	$GOVCONS$ para 48. En esta última variable, los 16 ceros consecutivos
	publicados para Venezuela entre 1996 y 2011 se tratan como faltantes por ser
	incompatibles con el dominio económico del indicador, sin sustituirlos ni
	imputarlos. Algunas economías con ECI disponible mediante API no cumplen los
	criterios del Atlas para aparecer en perfiles y rankings públicos; esta
	condición se documenta como limitación de calidad, pero no se utiliza para
	mezclar fuentes ni definir otra muestra econométrica. Las demás variables
	también presentan vacíos puntuales o años
	estructuralmente no publicados. Estas ausencias no se imputan. La muestra
	efectiva de cada una de las dos ecuaciones será la intersección de las
	variables requeridas y se documentará junto con las estimaciones.
	Aunque varios insumos llegan hasta 2022, las estimaciones cubrirán
	1996--2021 porque los cuatro componentes de $RENTS$ terminan en 2021.


	\section{Definición y construcción de variables}
	\label{sec:definicion-construccion-variables}

	Dado que las definiciones operativas completas de las variables del modelo ya fueron establecidas en el capítulo metodológico, esta sección se limita a precisar cuáles indicadores se utilizan en la descripción del panel y en el análisis exploratorio previo a la econometría. En esta etapa se privilegian las variables que permiten caracterizar tres dimensiones centrales de la base: la especialización extractiva, el nivel de desarrollo económico y la transformación estructural.

	Para la caracterización de la especialización extractiva se emplean indicadores de participación de las exportaciones de recursos naturales no renovables del subsuelo ---hidrocarburos, carbón, minerales, metales, piedras preciosas y oro no monetario--- en las exportaciones totales de mercancías, así como la suma de las rentas de petróleo, gas natural, carbón y minerales como porcentaje del PIB. Para la transformación estructural externa se utiliza principalmente el Índice de Complejidad Económica ($ECI$), complementado con $DIVX=1-HHI$. Por su parte, el nivel de desarrollo se resume mediante el logaritmo del PIB per cápita en PPA constante y, en algunos ejercicios descriptivos, mediante la clasificación de países por grupos de ingreso. Las transformaciones aplicadas respetarán exclusivamente las definiciones operativas del capítulo metodológico.

	En consecuencia, este capítulo no vuelve a desarrollar variable por variable la lógica del modelo econométrico. La capa \texttt{data/processed} reúne el panel depurado que se utilizará para documentar regularidades empíricas básicas antes de pasar a las estimaciones. Las definiciones detalladas, los criterios exactos de construcción y la justificación de las variables de control se mantienen en el capítulo de metodología para evitar redundancias innecesarias.


	\section{Estadística descriptiva}
	\label{sec:estadistica-descriptiva}

	Esta sección ofrecerá una caracterización descriptiva de la base empírica una vez construida la capa \texttt{data/processed}. El objetivo será identificar regularidades básicas en la distribución de las variables, en la composición de la muestra y en las asociaciones bivariadas más relevantes antes de proceder a la estimación econométrica. Por ahora se conservan únicamente gráficos de contextualización basados en fuentes externas claramente identificadas.

	\subsection{Hechos estilizados}
	\label{subsec:hechos-estilizados}

	Antes de examinar el panel consolidado de esta investigación, resulta útil presentar algunos hechos estilizados que sirven como punto de referencia para interpretar la heterogeneidad de las economías dependientes de recursos naturales no renovables del subsuelo. Los gráficos presentados en las Figuras \ref{fig:extractive-specialization}, \ref{fig:commodities-income}, \ref{fig:energy-income-boxplot} y \ref{fig:mining-income-boxplot} se construyen a partir de los datos reportados en las tablas de \textcite[27--30]{Anne2021}. Para estos gráficos, la participación extractiva se calcula como la suma de las exportaciones mineras y energéticas; por tanto, no incluye productos agrícolas ni otros recursos renovables. Su función en esta sección es contextual: permiten situar el problema empírico general antes de pasar al análisis exploratorio del panel utilizado en la tesis.

	Un primer hecho estilizado surge al analizar conjuntamente la participación de las exportaciones mineras y energéticas en la canasta exportadora de los países considerados. El diagrama de dispersión de la \autoref{fig:extractive-specialization} muestra que las observaciones no se distribuyen de manera homogénea en el plano, sino que tienden a agruparse en tres configuraciones. En una de ellas predominan las economías con alta participación de exportaciones energéticas y baja participación de exportaciones mineras; en otra se ubican las economías con especialización minera marcada y escasa presencia de recursos energéticos; y entre ambos extremos aparece un grupo de economías mixtas con participaciones moderadas en ambos tipos de recursos. En conjunto, el gráfico sugiere que la dependencia extractiva suele adoptar perfiles relativamente especializados, mientras que los casos de alta dependencia simultánea en minería y energía son menos frecuentes.

	\begin{figure}[H]
		\centering
		\includegraphics[width=\textwidth]{\subfix{../../figures/extractive_specialization_scatter}}

		\caption[Exportaciones mineras y energéticas]{Relación entre la participación de las exportaciones mineras y energéticas en la canasta exportadora de los países considerados. Las líneas punteadas indican umbrales del 20\% que permiten distinguir distintos patrones de especialización extractiva.}

		\label{fig:extractive-specialization}
		\caption*{\footnotesize \textit{Fuente: Elaboración propia con base en datos reportados en \textcite{Anne2021}.}}
	\end{figure}

	Un segundo hecho estilizado aparece al relacionar la dependencia exportadora de recursos extractivos con el nivel de ingreso de los países. La \autoref{fig:commodities-income} sugiere que la especialización en exportaciones mineras y energéticas no disminuye de manera monotónica a medida que aumenta el ingreso per cápita. Aunque las economías de bajos ingresos exhiben una elevada dispersión, los grupos de ingreso medio alto y alto también presentan niveles considerables de dependencia exportadora. En particular, la mediana de la participación extractiva se mantiene relativamente elevada en estos grupos, lo que indica que la especialización en recursos naturales no renovables del subsuelo puede coexistir con niveles altos de ingreso.

	\begin{figure}[H]
		\centering
		\includegraphics[width=\textwidth]{\subfix{../../figures/p_commodities_income_boxplot}}

		\caption[Dependencia exportadora extractiva por nivel de ingreso]{Distribución de la participación de las exportaciones mineras y energéticas en el total de exportaciones según grupo de ingreso: bajos ingresos (LIC), ingresos medios bajos (LMIC), ingresos medios altos (UMIC) y altos ingresos (HIC).}

		\label{fig:commodities-income}
		\caption*{\footnotesize \textit{Fuente: Elaboración propia con base en datos reportados en \textcite{Anne2021}.}}
	\end{figure}

	Un tercer hecho estilizado se observa al comparar la intensidad de la dependencia exportadora entre distintos tipos de recursos extractivos. Las Figuras \ref{fig:energy-income-boxplot} y \ref{fig:mining-income-boxplot} muestran que las economías especializadas en petróleo, gas y minería alcanzan patrones diferenciados de concentración exportadora. Mientras las exportaciones energéticas pueden representar una fracción dominante de la canasta exportadora, la participación de las exportaciones mineras tiende a ser menor y más dispersa entre grupos de ingreso. Este contraste sugiere que no todas las economías extractivas enfrentan el mismo tipo de restricción estructural: la especialización basada en petróleo, gas y minería puede asociarse con patrones distintos de concentración y transformación estructural.

	\begin{figure}[H]
		\centering

		\begin{subfigure}{0.49\textwidth}
			\centering
			\includegraphics[width=\textwidth]{\subfix{../../figures/p_energy_income_boxplot}}
			\caption{Dependencia de exportaciones energéticas por grupo de ingreso.}
			\label{fig:energy-income-boxplot}
		\end{subfigure}
		\hfill
		\begin{subfigure}{0.49\textwidth}
			\centering
			\includegraphics[width=\textwidth]{\subfix{../../figures/p_mining_income_boxplot}}
			\caption{Dependencia de exportaciones mineras por grupo de ingreso.}
			\label{fig:mining-income-boxplot}
		\end{subfigure}

		\caption[Dependencia exportadora por tipo de recurso]{Comparación de la participación de las exportaciones de recursos energéticos y mineros en el total de exportaciones de los países considerados.}

		\label{fig:energy-mining-income-comparison}
		\caption*{\footnotesize \textit{Fuente: Elaboración propia con base en datos reportados en \textcite{Anne2021}.}}

	\end{figure}

	\section{Análisis exploratorio de datos}
	\label{sec:analisis-exploratorio}

	El análisis exploratorio se incorporará a partir del panel ya construido y validado en \texttt{data/processed}. En particular, deberán regenerarse con la definición definitiva de $RENTS$ ---suma de las rentas de petróleo, gas natural, carbón y minerales--- y con la serie ECI completa obtenida mediante API todos los gráficos que relacionen rentas extractivas, complejidad y diversificación. Hasta entonces no se presentan asociaciones calculadas con el flujo anterior.

	\thesisendchap{datos-variables}

\end{document}
